% ECONOMICS_PROBLEM: {"id": "D72-449", "title": "Authoritarian resilience and economic institutions", "jel_code": "D72", "source_id": "OP-0128"}
% !TeX program = xelatex
\documentclass[11pt,a4paper]{article}
\usepackage[margin=23mm]{geometry}
\usepackage{fontspec}
\setmainfont{Latin Modern Roman}
\usepackage{amsmath,amssymb,mathtools}
\usepackage{xcolor,enumitem,booktabs,longtable,array,xurl,hyperref}
\definecolor{muted}{HTML}{53636C}
\hypersetup{colorlinks=true,urlcolor=blue,linkcolor=blue}
\setlength{\parindent}{0pt}
\setlength{\parskip}{5pt}
\newcommand{\R}{\mathbb R}
\newcommand{\E}{\mathbb E}
\newcommand{\N}{\mathbb N}
\newcommand{\ind}{\mathbf 1}
\newcommand{\Var}{\operatorname{Var}}
\newcommand{\Cov}{\operatorname{Cov}}
\newcommand{\supp}{\operatorname{supp}}
\DeclareMathOperator*{\argmin}{arg\,min}
\DeclareMathOperator*{\argmax}{arg\,max}
\newcommand{\entrymeta}[3]{{\sffamily\small\color{muted}\textbf{\texttt{#1}}\quad|\quad #2\par #3\par}}
\newcommand{\Problemref}[1]{\texttt{#1}}
\newcommand{\JELref}[2]{\texttt{#2} (JEL \texttt{#1})}
\begin{document}
\section*{Authoritarian resilience and economic institutions}
\textbf{Problem ID:} \texttt{D72-449}\quad \textbf{JEL:} \texttt{D72}\par
% BEGIN ECONOMICS STATEMENT
\label{problem:OP-0128}
\label{atlas:prob:Q8}
{\sffamily\small\color{muted}\textbf{Q8} | Political economy | Editorial research problem family\par}
\subsection*{Definitions and assumptions}
Fix all regimes entering a declared observation window, including failed regimes. Let $R_{ct}\in\{0,1\}$ indicate authoritarian survival, $Y^*_{ct}>0$ true output, $\widetilde Y_{ct}>0$ reported output, and $L_{ct}$ satellite night lights. Let $a$ be a feasible change in elite accountability or information institutions. A model specifies citizen/elite strategies and transitions, together with
\[
\log\widetilde Y_{ct}=\log Y^*_{ct}+b_{ct}+u_{ct},\qquad
L_{ct}=g(Y^*_{ct},S_{ct})+\eta_{ct},
\]
where $b_{ct}$ is reporting distortion, $S_{ct}$ includes sectoral composition and electrification, and the joint disturbance law is declared. Proxy relationships require validation or bounded departures; they are not assumed invariant across regimes by notation. Require finite first moments of $|\log Y^*_{c,t+h}(a)|$ under both compared interventions.
\subsection*{Formal research question}
Identify or sharply bound the survival contrast $\mathbb E[R_{c,t+h}(a)-R_{c,t+h}(a_0)]$, the true-growth contrast
\[
\theta_h=\mathbb E[\log Y^*_{c,t+h}(a)-\log Y^*_{c,t+h}(a_0)],
\]
and the distribution of $b_{ct}$. Determine what external measurement anchors, political variation and equilibrium restrictions distinguish productive accountability from manipulation and repression. For a selected sample of surviving regimes, characterize and address the additional survival-selection problem before transporting results to all entrants.
\subsection*{Interpretation and scope}
Night lights provide an alternative measurement system; they are not exact GDP observations. Simultaneously shifting reporting and the output-to-light mapping can leave both observed series unchanged, so identification needs restrictions or additional anchors. Growth performance and survival are distinct outcomes. The problem does not impose a normative objective of sustaining autocracy; it is a positive question about institutional mechanisms and their economic effects.
\subsection*{Source audit and status}
All three original citations are verified, with Besley--Kudamatsu resolved to the 2008 chapter and distinguished from its 2007 manuscript and an unrelated same-title book. Guriev--Treisman provides information-control mechanisms; Martinez evaluates reporting distortions. Neither proxy evidence nor selectorate theory establishes a universal causal growth effect. The displayed observation system and policy contrasts are editorial.
\subsection*{References}
{\footnotesize\raggedright
\begin{enumerate}[label={\textbf{[S\arabic*]}},leftmargin=3em]
\item Timothy Besley and Masayuki Kudamatsu (2008). \emph{Making Autocracy Work}. In Elhanan Helpman (ed.), Institutions and Economic Performance, Harvard University Press, 452--510.\par
\textit{Locator and source role:} Author publication list verifies chapter/date; model of selectorate accountability in author manuscript. Distinct from the unrelated 2016 book of the same title.\par
\url{https://sites.google.com/site/mkudamatsu/}
\item Sergei Guriev and Daniel Treisman (2019). \emph{Informational Autocrats}. Journal of Economic Perspectives 33(4), 100--127.\par
\textit{Locator and source role:} Abstract and discussion of information control; mechanisms of authoritarian rule, not causal identification of every regime transition.\par
\url{https://www.aeaweb.org/articles?id=10.1257/jep.33.4.100}
\item Luis R. Martinez (2022). \emph{How Much Should We Trust the Dictator's GDP Growth Estimates?}. Journal of Political Economy 130(10), 2731--2769.\par
\textit{Locator and source role:} Abstract and GDP/night-lights comparison; measurement of reporting distortions under maintained proxy assumptions.\par
\url{https://www.journals.uchicago.edu/doi/10.1086/720458}
\end{enumerate}
}

\subsection*{Common conventions: Source mathematical conventions}

Unless an entry states otherwise, agent and firm sets are finite. State and action spaces are measurable, strategies and outcome maps are measurable, probability laws are specified, and expectations exist for the targets under discussion. Infinite-horizon discount factors lie in $(0,1)$ and discounted objectives are finite. Norms, tolerances and complexity input sizes must be specified for computational comparisons. Constants or primitives that appear locally are inputs, not empirical facts asserted by this catalogue.

\textbf{Identification.} A statistical model $m\in\mathcal M$ implies an observable law $P_m$ and target $\theta(m)$. At law $P$, its identified set is
\[
\Theta(P;\mathcal M)=\{\theta(m):m\in\mathcal M,\ P_m=P\}.
\]
Point identification holds when this set is a singleton, or equivalently when $P_{m_1}=P_{m_2}$ implies $\theta(m_1)=\theta(m_2)$ for all compatible models. Sharp bounds mean bounds attained, or approached when the boundary is not attained, by compatible models. Writing this set is a definition, not a solution: the substantive task is to establish informative restrictions and derive or estimate the set. An unrestricted-confounding model may give only trivial bounds.

\textbf{Inference.} Identification, estimability and valid uncertainty quantification are separate questions. Fix $\alpha\in(0,1)$. A confidence procedure $C_n$ over a declared law class $\mathcal P$ has asymptotic uniform coverage at level $1-\alpha$ when
\[
\liminf_{n\to\infty}\inf_{P\in\mathcal P}\Pr_P\{\theta(P)\in C_n\}\geq1-\alpha.
\]
For set-identified targets the coverage event must be defined separately, for example coverage of the full identified set. If an entry uses identification-indexed models instead of uniquely indexed laws, the same coverage convention is applied over those models. Pointwise limits do not establish uniform validity, and asymptotic validity does not guarantee finite-sample precision. Sample sizes, dependence structures and weak-identification sequences are part of each application.

\textbf{Counterfactuals.} $Y_i(a)$ is a potential outcome under a specified intervention, including its timing, implementation and versions. It is not $Y_i$ conditional on voluntarily choosing $a$. A full intervention vector permits interference; shorthand scalar treatment notation does not silently impose no interference. Assignment, selection, attrition, anticipation and measurement must be specified. A claim about an instrument's local effect is not automatically a population policy effect.

\textbf{Economic equilibrium.} A counterfactual model includes best responses, constraints, feasibility, market clearing, state transitions and expectations consistent with the stated information. When several equilibria exist, either a declared selector is an input or the target is set-valued. Comparisons across policy regimes must use the stated selection convention. Reduced-form relationships are not presumed invariant after an equilibrium-changing intervention.

\textbf{Optimization and welfare.} Optimization is over the stated feasible set. If attainment is not established, $\sup$ or $\inf$ and $\varepsilon$-optimal choices replace an asserted maximizer or minimizer. For example, with value $V=\sup_{a\in\mathcal A}W(a)<\infty$, an $\varepsilon$-optimal set is $\{a\in\mathcal A:W(a)\geq V-\varepsilon\}$ for $\varepsilon>0$. Benefits, costs, risk and health or environmental outcomes can be aggregated only using declared units and conversion weights. Interpersonal utility weights, the treatment of fiscal transfers and foreign profits, discounting, and the behavioral welfare criterion are explicit inputs. A comparative-static derivative additionally requires existence and appropriate differentiability of the selected counterfactual.

\textbf{Sources.} Local labels [S1], [S2], and so forth restart in every question. Locators identify the relevant abstract, model, section or result at the level actually checked; a bibliographic or abstract-level verification is not represented as inspection of an entire proof. Working papers are distinguished from later publications and changing versions. Source summaries are paraphrased. All URL checks and editorial formulations refer to the audit date stated above.
% END ECONOMICS STATEMENT
\end{document}
